\section{Discussion}
\label{sec:discussion}

\subsection{Verdict on each clause}

Clause (i) is \textbf{refuted} against its own preregistered criterion: the pooled estimate is
$r = +0.026$ with a CI that includes zero, and equivalence testing rejects a true correlation as
large as $r = 0.20$. Clause (ii) is \textbf{refuted on both criteria}:
$\Delta\mathrm{SUVR} = +0.005$ per 10 \pp with a CI including zero, and a magnitude of 1.8\% of
the amyloid-negative/positive separation against a 5\% threshold. Clause (iii) is
\textbf{not supported and contradicted in direction}: all three paired contrasts are negative,
and the reliability null model predicts \stage should lead \tst by $1.86\times$ when the observed
ratio is $0.16\times$. Because the hypothesis is a conjunction, it fails.

\subsection{Stage is not spectrum, and the spectrum is one cohort}

The most useful result here is not the \stage null in isolation --- it is the dissociation
combined with the cohort concentration. The mechanistic story \citep{xie2013,fultz2019} is about
slow \emph{oscillations} driving CSF flow, and the human evidence supporting it is spectral:
$r = -0.45$ \citep{mander2015}, $-0.36$ \citep{winer2019}, $-0.52$ \citep{winer2020}. Every one
of those is the Berkeley Aging Cohort Study. Pooled across all three cohort families that have
measured spectral slow-wave activity, the estimate is $r = +0.112$ with $I^2 = 77\%$; with
Berkeley removed it is $-0.025$. Meanwhile the \emph{stage} metric --- the one clinicians read
off a sleep study, and the one the hypothesis names --- is essentially zero and homogeneously
so, with four independent cohort families agreeing.

The literature's habit of writing ``slow-wave sleep'' for both is therefore not a semantic
quibble. The construct with a positive but fragile signal and the construct with a tight null
are different measurements, and only one of them has ever produced the effect that motivates the
field's interest. We do not refute the refined, spectral version of the hypothesis: our spectral
stratum has a minimum detectable effect of $r = 0.417$ and rejects no equivalence bound, so its
honest verdict is insufficient evidence in either direction.

\subsection{How much to believe the total-sleep-time result}

\tst comes out ahead, which the hypothesis explicitly predicts should not happen, but three
things temper it. It is not itself significant at $\alpha = 0.05$ ($p = 0.052$, permutation
$p = 0.062$). It shows the clearest funnel asymmetry of any stratum, and trim-and-fill nudges it
to $r = 0.143$. And its $I^2$ is 77\%, driven by disagreement between the very large self-report
cohorts (\afour, $r \approx 0.02$) and the small PSG cohorts (\berkeley $0.36$, \blsa $0.33$) ---
itself a small-study signature. The fair statement is not that \tst beats \stage and is real, but
that \tst is the only sleep metric here with a defensible non-zero signal, and even that signal
is fragile, while \stage has no signal at all under any specification.

\subsection{What the reliability model adds}

Without E3, the finding ``\stage $\approx 0$ while \tst $\approx 0.16$'' would be ambiguous:
perhaps N3 is simply noisier. E3 shows the reverse, and disposes of a common defence of the
hypothesis --- that single-night polysomnography is too noisy to detect the N3 effect. It is too
noisy for total sleep time; it is not too noisy for N3. The literature is at its cleanest exactly
where it finds nothing. We designed this analysis expecting it to \emph{excuse} part of the
shortfall, and it became the strongest quantitative argument against clause (iii). We think the
same null model is worth running whenever a field claims one measure of a construct outperforms
another.

A second methodological lesson comes from the specification curve: $\tau^2$ estimator, $\rho$,
degrees-of-freedom rule and fixed-versus-random effects barely move anything, but whether
reported nulls and unstandardised coefficients are included moves the spectral estimate from
$+0.444$ to $+0.112$. In this literature, extraction completeness dominates statistical
technique.

\subsection{Limitations}

\para{Study-level design.} This is a meta-analysis of published effects, not individual
participant data. ADNI, A4, BLSA, AIBL, WRAP, \mayo, \agewell and OASIS-3 all require data-use
agreements with human review that could not be completed here. A genuine per-10\%-N3
dose-response \emph{curve} therefore cannot be estimated: \secref{sec:e4} converts pooled
standardised effects using external dispersion and reference scales, which assumes linearity
over the range and that the external SDs apply to the study populations.

\para{Power.} With 3--6 cohort families per stratum, $\tau^2$ is estimated unstably, the spectral
prediction interval spans the whole scale, and funnel-based diagnostics are near-powerless. HKSJ,
exact permutation and TOST keep the small-$k$ inference honest but cannot manufacture
information. The \stage stratum is precise enough to exclude $r \geq 0.20$ and no more.

\para{Conversion assumptions.} Test-statistic recovery assumes the reported $p$-value
corresponds to the coefficient's $t$-test with the stated degrees of freedom; covariate counts
were not fully recoverable from the text for 2 of 13 adjustment strings, though S3 shows the
answer is insensitive to a $\pm n_{\mathrm{cov}}$ change. Imputing $z = 0$ for narrative nulls
understates their true uncertainty; S1 removes them and the conclusion is unchanged. The
dependency correlations $\rho = 0.6$ and $\rho_w = 0.5$ are assumptions, not estimates; both
were varied without changing any conclusion.

\para{Reliability source.} The ICCs come from \sleepedf --- a healthy, mostly non-clinical
sample of 75 subjects on two consecutive nights, scored by the older Rechtschaffen--Kales
criteria with S3+S4 merged to N3. Consecutive-night ICCs include a first-night effect and may
understate the reliability of a habitual measure, and the amyloid cohorts are older and more
clinically heterogeneous. We therefore treat E3 as an order-of-magnitude prediction, and note
that the conclusion holds across the full CI of every ICC and under an alternative self-report
reliability assumption. Spectral slow-wave activity has no reliability estimate here at all,
because spectral power requires the raw EEG rather than the hypnogram annotations we hold, so
that stratum is not disattenuated.

\para{Corpus completeness.} Of 74 selected papers, 54 yielded full text; 17 of those were
recovered by typesetting PMC full-text HTML because PMC blocks automated PDF retrieval, so
figures and image-based tables are absent from them and a figure-only effect size could have
been missed. Twenty records remained abstract-only. All extractions are single-coder with
verbatim provenance quotes rather than dual independent extraction --- a real deviation from
PRISMA practice.

\para{Threats we cannot rule out.} Residual confounding by age, APOE4 status and apnea differs
across contributing studies and is only partly adjusted. Publication bias in the \emph{stage}
literature would most plausibly suppress positive N3 findings --- nulls in this area are usually
reported in passing, as ours largely were --- which would bias our estimate downward, though the
bias diagnostics give no sign of it. Finally, a real N3 effect confined to a specific brain
region or a specific sub-band would be invisible to a global-burden meta-analysis. Our null is
about global amyloid PET and stage-level N3, and it does not exclude regionally or spectrally
specific mechanisms.

\subsection{Broader implications}

\para{For trial design.} Targeting N3 stage minutes as a surrogate for amyloid clearance is not
supported by the aggregate human evidence. If slow-wave enhancement is to be tested, the exposure
should be defined spectrally (proportion of $<$1\,Hz activity), the trial should be powered for
effects well below $r = 0.45$, and it should not assume the Berkeley-scale effect will replicate.

\para{For measurement practice.} The stage and spectral definitions should be reported separately
and never collapsed. Studies should report both from the same recording, as \citet{winer2020}
did; that single practice is what made the decisive comparison in this analysis possible.

\para{For meta-analytic practice.} Restricting to natively-reported correlations gives a spectral
estimate of $+0.444$ rather than $+0.112$. Reported nulls and unstandardised coefficients are not
optional extras in this area --- they carry most of the corrective information, and any synthesis
that discards them will recover the literature's headline effect rather than test it.
